\section{Introduction}\label{sec:intro}

The $k$-nearest neighbor join (kNN join) computes, for every query object in a set $U$, its $k$ nearest neighbors in a reference set $I$. Unlike single-query kNN search, kNN join materializes a complete neighborhood relation between two datasets. This operation is a core primitive in similarity analytics, data mining, recommendation, classification, and anomaly detection, where large collections of objects are commonly represented as feature vectors~\cite{bohm2004k,xia2004gorder,yu2007efficient,hu2021efficient}.

Another defining characteristic of current similarity-analytics workloads is the growing prevalence of high-dimensional data representations. As dimensionality grows, full-distance computation becomes expensive and simple geometric pruning becomes less effective because nearest and farthest distances tend to become less distinguishable~\cite{beyer1999nearest,kouiroukidis2011effects,gao2023high,zhao2023towards}. High dimensionality therefore remains a central obstacle to efficient exact kNN join maintenance.

Many update-intensive applications require neighborhood relations to be maintained under continuous updates. In clinical databases, financial monitoring, and scientific or engineering pipelines, newly arriving cases or observations must be compared against evolving reference collections, while historical records may themselves be corrected, retired, or refreshed~\cite{an2023comprehensive,immadisetty2025real,li2024tracking,wei2016multisensor,giommi2023machine}. For example, in high-dimensional clinical case reanalysis, new patient cases must be matched against an evolving reference cohort while prior records may also be corrected or withdrawn, and an incorrect maintained neighborhood relation can directly mislead downstream diagnosis and treatment decisions. This motivates the fully dynamic kNN join problem: maintain kNN join results when both the query set $U$ and the reference set $I$ evolve.

The dual-update setting is substantially harder than the static or one-sided dynamic settings. An update to $U$ requires computing kNN results over $I$ to create or repair join rows. An update to $I$ requires identifying all queries in $U$ whose current kNN results are affected, and deletions from $I$ additionally require repairing those invalidated rows with fresh kNN searches. A useful way to view the task is as incremental maintenance of a materialized neighborhood view over two mutable inputs. Existing dynamic kNN join methods are mainly partially dynamic: they accelerate only one update direction or assume one side of the join is fixed. A single index on only one side of the join therefore accelerates only one side of this update process while leaving the other as a bottleneck.

Approximate nearest-neighbor methods are an important response to high-dimensional search and are well suited to many large-scale retrieval workloads~\cite{bawa2005lsh,jegou2010product,muja2014scalable,li2019approximate,gao2023high,zhao2023towards}. However, they optimize a different objective from the one studied in this paper. We focus on exact kNN join maintenance, where the materialized join table must remain identical to the result of recomputation under the chosen distance metric and tie-breaking rule. This exactness is practically important in real applications when the maintained join is used in high-stakes healthcare workflows, financial monitoring, clinical database reanalysis, or costly scientific and engineering pipelines, where stable and reproducible neighbor relations are valuable~\cite{an2023comprehensive,al2021financial,immadisetty2025real,li2024tracking,wei2016multisensor,giommi2023machine}. In the fully dynamic setting, approximation is particularly delicate because each update may involve both kNN-based row repair and affected-query identification; inaccuracies introduced in either stage can propagate through subsequent maintenance and cause the materialized join to drift from the exact result. Thus, the value of exactness here is not only that it can provide more faithful neighborhood semantics under the chosen metric, but also that it provides a stronger correctness guarantee for dynamic maintenance. This need is especially concrete in high-dimensional clinical case reanalysis, where evolving clinical databases and patient-similarity-based inference make stable maintained neighborhoods directly consequential for diagnosis and treatment decisions~\cite{li2024tracking,fang2021patient,gottlieb2013method}.

These observations lead to a natural maintenance view of the problem. Fully dynamic kNN join can be understood as two coordinated maintenance paths over a materialized neighborhood view: one computes or repairs rows through kNN search after updates to $U$, whereas the other first identifies affected queries and then repairs only those rows after updates to $I$. To handle these two paths in a coordinated manner, we propose a bidirectional locate-and-repair framework.

To realize this framework in the exact high-dimensional setting, the underlying search structures must still support effective reduced-space pruning, exact leaf-level verification, and repeated maintenance. We address this requirement with \dualtree, which instantiates the framework using a \deltatree on the reference set $I$ for kNN search and row repair and a dynamically maintained \hdrtree on the query set $U$ for affected-query localization, while keeping both indexes synchronized with the materialized join result.

On top of this framework instantiation, we introduce a point-wise low-dimensional pruning stage at leaf nodes, pushing pruning beyond prior cluster-level reduced-space checks. More importantly, we derive a cost-aware analytical model for selecting the pruning dimensionality used by this fine-grained pruning stage. The model relates the expected computation cost to the pruning probability as a function of dimension and yields an optimal or near-optimal dimension for the leaf-level check, replacing heuristic dimension selection with an explicit optimization objective.

In summary, this paper makes the following contributions:

\noindent(1) \emph{A bidirectional locate-and-repair framework.}
We formulate exact fully dynamic kNN join maintenance as a bidirectional locate-and-repair problem over a materialized neighborhood view. This abstraction separates kNN-based row repair from affected-query localization and provides the framework-level basis of our method.

\noindent(2) \emph{A dual-index instantiation.}
We instantiate this framework with \dualtree, which combines \deltatree-based kNN search and row repair with a dynamically maintained \hdrtree-based affected-query localization component. The resulting system supports updates on both the query and reference sides while keeping the materialized join table and index summaries synchronized throughout maintenance.

\noindent(3) \emph{Point-wise pruning and cost-aware modeling.}
We extend reduced-space pruning from the cluster level to the point level at leaf nodes and derive a cost-aware model that selects the projection dimensionality by balancing pruning power against projection cost. This replaces heuristic dimension selection with an explicit optimization objective.

\noindent(4) \emph{Comprehensive evaluation.}
We evaluate the proposed methods on eight real datasets with dimensionalities from 128 to 4096. The results show that \dualtree consistently outperforms strong one-sided baselines, while \dualtreeopt achieves an average speedup of 1.88$\times$ over the strongest one-sided baseline across seven additional datasets and up to 2.65$\times$ on Trevi.

The remainder of the paper is organized as follows. Section~\ref{background} gives definitions and index preliminaries. Section~\ref{sec:related-works} reviews related work. Section~\ref{sec:framework} presents the bidirectional framework and its dual-index instantiation. Section~\ref{sec:pruning} develops fine-grained pruning and cost-aware dimension selection. Section~\ref{sec:exp} reports the experimental results, and Section~\ref{sec:conclusion} concludes the paper.
